\chapter{STRATEGIC MARKET ANALYSIS}
\label{ch:market}

\section{The market segment}

\subsection{The potential market}

The service addresses anyone who can lose a personal object and reach a web
platform to say so. In Algeria that population is large and, more importantly,
almost entirely connected.

\begin{table}[htbp]
\centering
\caption{Addressable population in Algeria}
\label{tab:market-size}
\small
\begin{tabularx}{\textwidth}{@{}X r l@{}}
\toprule
\textbf{Indicator} & \textbf{Value} & \textbf{Source} \\
\midrule
Population (Q2 2025 estimate)            & 47.40 M  & ARPCE~\cite{arpce2025} \\
Internet users (January 2025)            & 36.20 M  & DataReportal~\cite{datareportal2025} \\
Internet penetration                     & 76.9\,\% & DataReportal~\cite{datareportal2025} \\
Internet subscriptions (Q2 2025)         & 59.10 M  & ARPCE~\cite{arpce2025} \\
\quad of which mobile                    & 88.71\,\% & ARPCE~\cite{arpce2025} \\
\quad of which fixed-line                & 11.29\,\% & ARPCE~\cite{arpce2025} \\
Mobile lines                             & 54.80 M  & ARPCE~\cite{arpce2025} \\
Mobile penetration rate                  & 115.77\,\% & ARPCE~\cite{arpce2025} \\
Urban population                         & $\approx$ 66.8\,\% & Macrotrends~\cite{macrotrends2025} \\
University students (2024--2025)         & $>$ 1.80 M & MESRS~\cite{mesrs2025} \\
Higher-education institutions            & $>$ 130   & MESRS~\cite{mesrs2025} \\
\bottomrule
\end{tabularx}
\end{table}

Three features of this market shape the product directly.

\paragraph{It is mobile-first, decisively.} Nearly nine internet subscriptions
in ten are mobile~\cite{arpce2025}. A citizen who has just lost a wallet is
standing in the street with a telephone, not sitting at a desk. This is why the
platform is a responsive web application requiring no installation rather than a
store-distributed app: at the moment of loss, the friction of downloading
anything is the friction that loses the report.

\paragraph{It is urban.} Roughly two thirds of Algerians live in urban
areas~\cite{macrotrends2025}, and loss events concentrate further still --- in
transport, markets, campuses and commercial centres. Density matters more than
national totals for this product, because matching only works where both a loss
and a discovery occur within the same place and period.

\paragraph{It is multilingual.} Arabic, French and English circulate together,
frequently inside a single sentence. Any system that compares reports by
keyword fragments this market into non-communicating sub-pools. The multilingual
encoder described in Chapter~\ref{ch:prototype} is a response to this market
structure, not a technical preference.

\subsection{The target market}

The potential market is not the market to launch into. Matching produces value
only at density: one thousand reports concentrated in a single wilaya are worth
far more than ten thousand spread across the country. The go-to-market therefore
narrows deliberately, as shown in Figure~\ref{fig:market-funnel}.

\begin{figure}[htbp]
\centering
\begin{tikzpicture}[x=1cm,y=1cm]

  \node[draw=uamobblue, fill=uamobblue!8, rounded corners=2pt, align=center,
        font=\footnotesize, text width=12.4cm, inner sep=6pt] (tam) at (0,4)
    {\textbf{TAM --- Total addressable market}\\
     36.2 million internet users in Algeria~\cite{datareportal2025}};

  \node[draw=uamobblue, fill=uamobblue!16, rounded corners=2pt, align=center,
        font=\footnotesize, text width=9.4cm, inner sep=6pt] (sam) at (0,2.2)
    {\textbf{SAM --- Serviceable available market}\\
     urban connected population, $\approx$ 24 million\\
     {\scriptsize(internet users $\times$ 66.8\,\% urban share)}};

  \node[draw=uamobblue, fill=uamobblue!30, rounded corners=2pt, align=center,
        font=\footnotesize, text width=6.4cm, inner sep=6pt] (som) at (0,0.4)
    {\textbf{SOM --- Serviceable obtainable market}\\
     pilot wilayas and partner campuses\\
     {\scriptsize(year one: Bouira, then Algiers)}};

  \draw[flowarrow] (tam.south) -- (sam.north);
  \draw[flowarrow] (sam.south) -- (som.north);

\end{tikzpicture}
\caption{Market funnel. The serviceable available market is an approximation
obtained by applying the national urban share to the connected population; it is
indicative rather than surveyed.}
\label{fig:market-funnel}
\end{figure}

The launch segment is therefore \textbf{connected urban residents of dense
wilayas, reached first through institutions rather than individually}. Three
characteristics define it:

\begin{itemize}[leftmargin=1.4em,itemsep=3pt]
  \item \textbf{Students and campus staff.} More than 1.8 million students study
        across over 130 institutions~\cite{mesrs2025}. A campus is a bounded
        space with high object turnover, an existing lost-property desk, and a
        population already comfortable with digital services --- the ideal pilot
        ground, and the reason the first deployment targets Akli Mohand Oulhadj
        University of Bouira.
  \item \textbf{Users of public transport and commercial centres}, where loss
        events cluster and where operators already hold unsearchable registers.
  \item \textbf{People who already report losses online}, in social-media
        groups. They have demonstrated the behaviour the platform needs; what
        they lack is a system that does the comparing.
\end{itemize}

\section{Competitor analysis}

\subsection{Direct competitors}

\paragraph{Social-media groups.} Dedicated Facebook groups for lost and found
objects exist in most Algerian cities and are, in practice, the incumbent. Their
strength is real: they are free, already populated, and require nothing new from
the user. Their limitation is structural rather than incidental. A post is
unstructured text in a reverse-chronological feed; it is not indexed, cannot be
filtered by category, place or date, and --- decisively --- nobody rereads last
month's posts against today's. The comparison work is left entirely to human
attention, which does not scale and does not persist.

\paragraph{Classified-ads platforms.} General marketplaces such as Ouedkniss
carry occasional lost-and-found listings. They provide structure and search, but
they are built for a seller who wants to be found by a buyer who knows what they
are looking for. Lost property inverts that: neither party knows the other
exists, and neither can formulate the query that would surface the other.

\paragraph{Institutional lost-property desks.} Universities, transport
operators and municipalities each hold objects and each keeps a register. These
are not really competitors --- they are the intended partners of
Chapter~\ref{ch:organization}. But as things stand they compete for the same
event, and they fail at it: the registers are private, local, paper-based, and
searchable only by physically visiting during opening hours.

\subsection{Indirect competitors}

\paragraph{Bluetooth trackers.} Devices such as Apple AirTag or Samsung
SmartTag locate an object that has been prepared in advance. They work well for
what they cover and cannot cover most of the problem: the tracker must be bought
and attached \emph{before} the loss, one per object, and it does nothing for a
wallet, a set of identity papers, a document folder or a piece of jewellery.
They also locate an object without connecting the finder to the owner.

\paragraph{Device-level recovery features.} \dq{Find My Device} functionality is
built into modern phones and is the correct tool for a lost phone. It applies to
exactly one category of object, and only while the device retains power and
connectivity.

\paragraph{Insurance and replacement.} Replacing what was lost is the substitute
of last resort. It is slow for administrative documents, expensive for
electronics, and impossible for anything of sentimental value.

\paragraph{Doing nothing.} The realistic default, and the true incumbent. Most
losses are never reported anywhere, because the perceived probability of
recovery does not justify the effort. This is the behaviour the product must
change, and it sets the standard the interface has to meet: reporting must take
less effort than giving up.

\begin{table}[htbp]
\centering
\caption{Strengths and weaknesses of the competition}
\label{tab:competitors}
\small
\begin{tabularx}{\textwidth}{@{}p{3.1cm} X X@{}}
\toprule
\textbf{Competitor} & \textbf{Strengths} & \textbf{Weaknesses} \\
\midrule
Social-media groups
  & Free; already populated; no new habit to learn; strong local reach
  & No structure, no index, no filtering; posts scroll out of sight in hours;
    no comparison of old reports against new; contact details exposed publicly \\
\addlinespace
Institutional desks
  & Physically hold the objects; trusted; already exist
  & One private register per desk; local only; in-person and office hours;
    no cross-institution search \\
\addlinespace
Classified-ads platforms
  & Structured listings; search; national audience
  & Built for a known-item query; neither party can formulate the search;
    lost property is a marginal use case with no dedicated workflow \\
\addlinespace
Bluetooth trackers
  & Reliable for a prepared object; precise location
  & Must be purchased and attached before the loss; one device per object;
    useless for papers, wallets, keyrings without a tag; does not connect
    finder to owner \\
\addlinespace
Doing nothing
  & Costs no effort
  & Recovers nothing \\
\bottomrule
\end{tabularx}
\end{table}

\subsection{Positioning}

Two axes separate these options, and together they leave a quadrant empty.
The first is \textbf{who performs the search} --- the user, or the system. The
second is whether the solution \textbf{works on any object without preparation
before the loss}. Figure~\ref{fig:positioning} places the alternatives.

\begin{figure}[htbp]
\centering
\begin{tikzpicture}[x=1cm,y=1cm]

  \fill[uamobblue!7, rounded corners=2pt] (0.15,0.15) rectangle (5.9,3.9);

  \draw[->, thick, draw=rulegray] (0,0) -- (6.2,0)
    node[right, font=\scriptsize, align=left] {search performed\\by the system};
  \draw[->, thick, draw=rulegray] (0,0) -- (0,4.2)
    node[above, font=\scriptsize, align=center] {works on any object,\\no preparation};

  \node[font=\scriptsize, anchor=north east] at (6.2,-0.15) {automatic};
  \node[font=\scriptsize, anchor=north west] at (0,-0.15) {manual};

  \node[draw=rulegray, fill=white, rounded corners=2pt, font=\tiny,
        align=center, inner sep=3pt] at (1.35,3.15)
    {Social-media\\groups};
  \node[draw=rulegray, fill=white, rounded corners=2pt, font=\tiny,
        align=center, inner sep=3pt] at (1.35,2.05)
    {Institutional\\desks};
  \node[draw=rulegray, fill=white, rounded corners=2pt, font=\tiny,
        align=center, inner sep=3pt] at (2.5,1.15)
    {Classified\\ads};
  \node[draw=rulegray, fill=white, rounded corners=2pt, font=\tiny,
        align=center, inner sep=3pt] at (4.75,0.75)
    {Bluetooth\\trackers};
  \node[draw=rulegray, fill=white, rounded corners=2pt, font=\tiny,
        align=center, inner sep=3pt] at (4.75,0.05)
    {\dq{Find My}\\device};

  \node[draw=uamobblue, fill=uamobblue!25, thick, rounded corners=2pt,
        font=\scriptsize\bfseries, align=center, inner sep=4pt] at (4.7,3.2)
    {\projectname};

\end{tikzpicture}
\caption{Competitive positioning. The upper-right quadrant --- automatic search
that works on any object without prior preparation --- is unoccupied.}
\label{fig:positioning}
\end{figure}

\section{SWOT analysis}

\subsection{Purpose of the SWOT analysis}

The analysis serves two purposes here. It states plainly what the project can
rely on and what it cannot, and it forces the weaknesses to be named before a
jury or an investor names them. The most consequential entry below is not a
strength: it is the cold-start problem, which determines whether the platform
works at all in its first year, and which the go-to-market of
Section~\ref{sec:marketing} is built specifically to answer.

\subsection{SWOT matrix}

\begin{table}[htbp]
\centering
\caption{SWOT matrix}
\label{tab:swot}
\small
\renewcommand{\arraystretch}{1.25}
\begin{tabularx}{\textwidth}{@{}>{\bfseries}p{2.3cm} X@{}}
\toprule
Strengths &
\begin{itemize}[leftmargin=1em,nosep,topsep=2pt]
  \item Multimodal matching engine combining semantic text and image similarity,
        already implemented and working
  \item Genuinely multilingual --- the same pool serves Arabic, French and
        English reports without fragmenting
  \item Privacy-preserving claim workflow: a match never exposes contact details
  \item No per-request AI cost; inference runs locally on open models, which is
        what makes free access sustainable
  \item Explained confidence rather than opaque ranking
  \item Built on open-source infrastructure with no licence fees
\end{itemize} \\
\midrule
Weaknesses &
\begin{itemize}[leftmargin=1em,nosep,topsep=2pt]
  \item \textbf{Cold start.} The platform is worth little until both pools are
        populated, and the found side is the scarcer one
  \item No revenue from the primary user, by design --- monetisation is indirect
        and unproven
  \item Scoring constants are hand-tuned; calibration awaits real feedback data
  \item Small team with no operational track record
  \item Dependent on voluntary honesty; verification questions mitigate but do
        not eliminate fraudulent claims
  \item Matching quality has been validated on a demonstration corpus, not in
        the field
\end{itemize} \\
\midrule
Opportunities &
\begin{itemize}[leftmargin=1em,nosep,topsep=2pt]
  \item 76.9\,\% internet penetration and a 115.77\,\% mobile penetration
        rate~\cite{datareportal2025,arpce2025} --- the audience is already
        connected
  \item No incumbent occupies the automatic-matching quadrant
        (Figure~\ref{fig:positioning})
  \item Institutions already hold unsearchable registers and gain from
        publishing them
  \item National Startup Label and ANADE financing mechanisms actively support
        ventures of this kind~\cite{startupdz,anade2025}
  \item Over 1.8 million students~\cite{mesrs2025} concentrated in bounded
        campuses --- dense, reachable pilot grounds
  \item The same engine extends to adjacent categories (documents, vehicle
        papers, pets) with no architectural change
\end{itemize} \\
\midrule
Threats &
\begin{itemize}[leftmargin=1em,nosep,topsep=2pt]
  \item A social-media platform adding structured lost-and-found features
  \item Abuse: fraudulent claims, spam reports, or use of the platform to locate
        people rather than objects
  \item Personal-data obligations around photographs and identity information
  \item Network effects favour whoever reaches density first --- being second is
        close to being absent
  \item Infrastructure costs scaling with image storage before revenue does
  \item Low trust in online handovers between strangers
\end{itemize} \\
\bottomrule
\end{tabularx}
\end{table}

\section{The marketing strategy}
\label{sec:marketing}

\subsection{Marketing strategy for the company}

The distinguishing feature of this market is that \textbf{demand is
event-driven}. Nobody seeks a lost-property service until the moment they lose
something, and at that moment they search rather than remember. Sustained
brand advertising is therefore poorly matched to the product; being present in
the search result, and being already populated when the user arrives, is what
converts.

The strategy follows from that, and from the cold-start weakness named above.

\begin{enumerate}[leftmargin=1.6em,itemsep=4pt]
  \item \textbf{Seed supply before demand.} Sign institutional partners first
        and publish their existing registers into the pool. A user arriving
        after a loss must find something to match against; an empty platform
        converts nobody and is never revisited.
  \item \textbf{Concentrate geographically.} Saturate one campus, then one
        wilaya, before widening. Match rate depends on density, and a thin
        national presence produces no matches anywhere.
  \item \textbf{Own the moment of search.} Optimise for the queries people
        actually type after a loss, in all three languages.
  \item \textbf{Be present where losses happen.} Physical signage at
        lost-property desks, transport hubs and campus buildings, carrying a
        direct link --- the cheapest acquisition channel available, placed
        exactly where the trigger occurs.
  \item \textbf{Lead with recovered objects.} Each successful recovery is the
        only marketing asset that convinces a sceptic, and it is free to
        produce.
\end{enumerate}

\subsection{The marketing mix (7P)}

\begin{itemize}[leftmargin=1.4em,itemsep=4pt]
  \item \textbf{Product.} A trilingual web platform that automatically compares
        every loss report against every discovery report, on text and images at
        once, and protects the handover behind a verification handshake.

  \item \textbf{Price.} Free for citizens, on both sides, permanently. This is a
        strategic constraint rather than a promotional one: charging the person
        who just lost their identity papers would suppress precisely the reports
        the engine needs, and charging the finder would suppress the scarcer
        side of the market outright. Revenue comes from institutions and
        partnerships instead --- see Chapter~\ref{ch:finance}.

  \item \textbf{Place.} The web, reachable from any smartphone browser without
        installation; and physically, the partner institutions where objects are
        handed in.

  \item \textbf{Promotion.} Search visibility in Arabic, French and English;
        signage at lost-property desks and transport hubs; campus partnerships;
        recovered-object stories; institutional press.

  \item \textbf{People.} A small technical team, with moderation and support
        capacity growing alongside claim volume --- disputes are handled by
        people, not by the algorithm.

  \item \textbf{Process.} Report, automatic matching, notification, verified
        claim, approved handover, closure. Each step is instrumented, and the
        outcome of each claim feeds the learning loop.

  \item \textbf{Physical evidence.} The interface itself; notifications;
        published recovery statistics; and the visible presence of the service
        at partner desks, which is what transfers institutional trust to a new
        platform.
\end{itemize}
